\RequirePackage[thinlines]{easybmat} %-- muss aufgrund eines Bugs vor etex und tikz geladen werden
\documentclass[a4paper, twoside, headsepline, index=totoc,toc=listof, fontsize=10pt, cleardoublepage=empty, headinclude, DIV=13, BCOR=13mm, titlepage,]{scrartcl}
\usepackage{scrtime} % KOMA, Uhrzeit ermoeglicht
\usepackage{scrpage2} % wie fancyhdr, nur optimiert auf KOMA-Skript, leicht andere Syntax
\usepackage{etoolbox}
\usepackage{letltxmacro}

% Compiler abhaengige IFs
\usepackage{ifxetex,ifluatex}
\newif\ifxetexorluatex
\ifxetex
  \xetexorluatextrue
\else
  \ifluatex
    \xetexorluatextrue
  \else
    \xetexorluatexfalse
  \fi
\fi

%--Farbdefinitionen
%-- muss vor tikz geladen werden
\usepackage[usenames, table, x11names]{xcolor}
\definecolor{dark_gray}{gray}{0.45}
\definecolor{light_gray}{gray}{0.6}

%--Zum Zeichnen
%-- muss vor polyglossia bzw. babel geladen werden
\usepackage{tikz}
\usepackage{tikz-cd}
\usetikzlibrary{external}
\tikzset{>=latex}
\usetikzlibrary{shapes,arrows,intersections}
\usetikzlibrary{calc,3d}
\usetikzlibrary{decorations.pathreplacing,decorations.markings}
\usetikzlibrary{angles}

%-- Konfiguration von tikzexternalize
\tikzexternalize[prefix=tikz/, up to date check=diff]

\ifxetexorluatex
\tikzset{external/system call={xelatex \tikzexternalcheckshellescape %-- verwende LuaLaTeX, wegen dynamischer Speicherallokation
    -halt-on-error -interaction=batchmode -jobname "\image" "\texsource"}}
\else
\tikzset{external/system call={pdflatex \tikzexternalcheckshellescape 
    -halt-on-error -interaction=batchmode -jobname "\image" "\texsource"}}
\fi

%-- tikzexternalize fuer tikzcd deaktivieren, da inkompatibel
\AtBeginEnvironment{tikzcd}{\tikzexternaldisable}
\AtEndEnvironment{tikzcd}{\tikzexternalenable}

%-- um Inkompatibilitaeten von quotes und polyglossia bzw. babel zu vermeiden
\tikzset{
  every picture/.append style={
    execute at begin picture={\shorthandoff{"}},
    execute at end picture={\shorthandon{"}}
  }
}
\usetikzlibrary{quotes}
\usepackage{pgfplots}
\usepgfplotslibrary{colormaps}
\pgfplotsset{compat=1.10}



%-- Mathe
\usepackage{mathtools} % beinhaltet amsmath
\mathtoolsset{showonlyrefs, centercolon}
\newtagform{brackets}[\textbf]{[}{]}
\usetagform{brackets}
\usepackage{wasysym}
\usepackage{amssymb} %zusätzliche Symbole
\usepackage{latexsym} %zusätzliche Symbole
\usepackage{stmaryrd} %für Blitz
\usepackage{nicefrac} %schräge Brüche
\usepackage{cancel} %Befehle zum Durchstreichen
\usepackage{mathdots}
% \usepackage{bm}
\DeclareSymbolFont{bbold}{U}{bbold}{m}{n}
\DeclareSymbolFontAlphabet{\mathbbold}{bbold}
\newcommand{\mathds}[1]{\mathbb{#1}} %Um Kompatibilitaet mit frueheren Benutzung von dsfont herzustellen
\newcommand{\ind}{\mathbbold{1}} % charakteristische-Funktion-Eins
\def\mathul#1#2{\color{#1}\underline{{\color{black}#2}}\color{black}} %farbiges Untersteichen im Mathe-Modus

%-- Alles was mit Schrift und XeteX zu tun hat
\ifxetexorluatex
    % XeLaTeX or LuaTeX
	\ifxetex
		\usepackage{mathspec} %beinhaltet fontspec 
	\else
		\usepackage[no-math]{fontspec}
	\fi
    \usepackage{polyglossia} %babel-ersatz
    \setmainlanguage[spelling=new,babelshorthands=false]{german}
	\newcommand\glqq{"}
	\newcommand\grqq{"}
	\defaultfontfeatures{Mapping=tex-text, WordSpace={1.4}} %
    \setmainfont[Ligatures=Common, BoldFont={* Bold}, ItalicFont={* Light Italic},SmallCapsFont={Linux Libertine O}, SmallCapsFeatures={Letters=SmallCaps}]{Source Sans Pro}
	\setsansfont[Scale=MatchLowercase,Ligatures=Common, BoldFont={* Medium}]{Ubuntu}
	\ifxetex
		\setallmonofonts[Scale=MatchLowercase, ItalicFont={*}]{Consolas} 
	\else
		\setmonofont[Scale=MatchLowercase, ItalicFont={*}]{Consolas}
	\fi
	\usepackage{xltxtra}
	\usepackage{fontawesome}
	\usepackage{microtype}
\else
    % default: pdfLaTeX
    \usepackage[ngerman]{babel}
    \usepackage[T1]{fontenc}
    \usepackage[utf8]{inputenc}
	\usepackage{textcomp} %verhindert ein paar Fehler bei den Fonts
    \usepackage[babel=true, tracking=true]{microtype}
	\usepackage{lmodern}
	\renewcommand{\familydefault}{\sfdefault} 
\fi


%--Mixed
\usepackage[neverdecrease]{paralist}
\usepackage[german=quotes]{csquotes}
\usepackage{makeidx}
\usepackage{booktabs}
\usepackage{wrapfig}
\usepackage{float}
\usepackage[margin=10pt, font=small, labelfont={sf, bf}, format=plain, indention=1em]{caption}
\captionsetup[wrapfigure]{name=Abb. }
\usepackage{stackrel}
\usepackage{ifthen}
\usepackage{multicol}
\usepackage{qrcode}
\flushbottom

%--Hyphenation
\hyphenation{di-men-si-o-na-le}


%--Unterstreichung
\usepackage[normalem]{ulem}
\setlength{\ULdepth}{1.8pt}

%--Indexverarbeitung
\newcommand{\bet}[1]{\textbf{#1}}
\newcommand{\Index}[1]{\textbf{#1}\index{#1}}
\makeindex
\setindexpreamble{{\noindent \itshape Die \emph{Seitenzahlen} sind mit Hyperlinks zu den entsprechenden Seiten versehen, also anklickbar} \ifxetexorluatex \faHandUp\fi \par \bigskip}
\renewcommand{\indexpagestyle}{scrheadings}


%--Marginnote & todonotes
\usepackage{marginnote}
\renewcommand*{\marginfont}{\color{dark_gray} \itshape \footnotesize }
\usepackage[textsize=small]{todonotes}

%--Todonotes mit tikz/externalize kompatibel machen
\LetLtxMacro{\oldtodo}{\todo}
\renewcommand{\todo}[2][]{\tikzexternaldisable\oldtodo[#1]{#2}\tikzexternalenable}
\LetLtxMacro{\oldmissingfigure}{\missingfigure}
\renewcommand{\missingfigure}[2][]{\tikzexternaldisable\oldmissingfigure[{#1}]{#2}\tikzexternalenable}

%--Konfiguration von Hyperref pdfstartview=FitH, 
\usepackage[hidelinks, pdfpagelabels,  bookmarksopen=true, bookmarksnumbered=true, linkcolor=black, urlcolor=SkyBlue2, plainpages=false,pagebackref, citecolor=black, hypertexnames=true, pdfauthor={Jannes Bantje}, pdfborderstyle={/S/U}, linkbordercolor=SkyBlue2, colorlinks=false]{hyperref}

% \let\orighref\href
\ifxetexorluatex
\newcommand{\hrefsym}[2]{\href{#1}{\texttt{#2\,\faExternalLink}}}
\newcommand{\hrefsymmail}[2]{\href{#1}{\texttt{\faEnvelope\,#2}}}
\renewcommand{\url}[1]{\hrefsym{#1}{\nolinkurl{#1}}}
\fi
\providecommand{\hrefsym}{\href}
% \usepackage{cleveref}
% \crefformat{footnote}{#2\footnotemark[#1]#3}

%--Römische Zahlen
\newcommand{\RM}[1]{\MakeUppercase{\romannumeral #1{}}}

%--Punkte (nach hyperref)
\usepackage{ellipsis}

%%--Abkürzungen etc.
\newcommand{\light}{\text{\Large $\lightning$}}

%-- Definitionen von weiteren Mathe-Befehlen
\DeclareMathOperator{\re}{Re} %Realteil
\DeclareMathOperator{\im}{Im} %Imaginaerteil
\DeclareMathOperator{\id}{id} %identische Abbildung
\DeclareMathOperator{\Sp}{Sp} %Spur
\DeclareMathOperator{\sgn}{sgn} %Signum
\DeclareMathOperator{\alt}{Alt} %Alternierende n-linearFormen
\DeclareMathOperator{\End}{End} %Endomorphismen
\DeclareMathOperator{\Vol}{Vol} %Volumen
\DeclareMathOperator{\dom}{dom} %Domain
\DeclareMathOperator{\Hom}{Hom} %Homomorphismen
\DeclareMathOperator{\bild}{Bild} %Bild
\DeclareMathOperator{\Kern}{Kern }
\DeclareMathOperator{\rg}{rg} %Rang
\DeclareMathOperator{\Rg}{Rg} %Rang
\DeclareMathOperator{\diam}{diam} %Durchmesser
\DeclareMathOperator{\dist}{dist} %Distanz
\DeclareMathOperator{\grad}{grad} %Gradient
\DeclareMathOperator{\dive}{div} %Gradient
\DeclareMathOperator{\rot}{rot} %Rotation
\DeclareMathOperator{\hess}{Hess} %Hesse-Matrix
\DeclareMathOperator{\koker}{Koker} %Kokern
\DeclareMathOperator{\aut}{Aut} %Automorphismen
\DeclareMathOperator{\ord}{ord} %Ordnung
\DeclareMathOperator{\ggT}{ggT} %ggT
\DeclareMathOperator{\kgV}{kgV} %kgV
\DeclareMathOperator{\Gr}{Gr} %Gerade
\DeclareMathOperator{\Kr}{Kr} %Kreis
\DeclareMathOperator{\Char}{char} %Charakteristik
\DeclareMathOperator{\Aut}{Aut} %Automorphismen
% \DeclareMathOperator{\D}{D} %formale Ableitung
\newcommand{\D}{\ensuremath{\mathrm{D}\mkern-1.5mu}}
\newcommand{\mathd}{\ensuremath{\mathrm{d}\mkern-1.0mu}}
\newcommand{\Tmap}{\ensuremath{\mathrm{T}\mkern-1.0mu}}
\newcommand{\opL}{\ensuremath{\mathrm{L}\mkern-0.6mu}}
\DeclareMathOperator{\cov}{cov} %Kovarianz
\DeclareMathOperator{\Gal}{Gal} %Galois
\DeclareMathOperator{\supp}{supp} %Träger
\DeclareMathOperator{\Sym}{Sym} %Symmetrische Gruppe
\DeclareMathOperator{\Zyl}{Zyl} %Zylinder
\DeclareMathOperator{\Mod}{Mod} %Moduln
\DeclareMathOperator{\EPK}{EPK} %Einpunktkompaktifizierung
\DeclareMathOperator{\conj}{conj} %Konjugation
\DeclareMathOperator{\ann}{ann} %Annulator
\DeclareMathOperator{\tor}{tor} %Torsionsmodul
\DeclareMathOperator{\Int}{Int} %Interior/Inneres
\DeclareMathOperator{\Br}{Br} %Brauerergruppe
\DeclareMathOperator{\GL}{GL} %allgemeine lineare Gruppe
\DeclareMathOperator{\SL}{SL} %spezielle lineare Gruppe
\DeclareMathOperator{\SU}{SU} %spezielle unitäre Gruppe
\DeclareMathOperator{\Tr}{Tr} %Spur/Trace
\DeclareMathOperator{\conv}{conv} %Konvexe Teilmenge
% \DeclareMathOperator{\tr}{tr}
\newcommand{\tr}{\mathrm{tr}}
\newcommand{\w}{\mathrm{w}}
\newcommand{\weakT}[1]{\ensuremath{\mathcal{T}_{#1}^{\mkern+1.0mu\text{\raisebox{0.4ex}{$\mathrm{w}$}}}}}
\newcommand{\weakTstar}[1]{\ensuremath{\mathcal{T}_{#1}^{\mkern+1.0mu\text{\raisebox{0.4ex}{$\mathrm{w}$}}^*}}}
\newcommand{\rand}[1]{\ensuremath{\partial^{\scriptscriptstyle #1}}}


% -- Kategorien
\DeclareMathOperator{\Ob}{Ob} %Objekte
\DeclareMathOperator{\Mor}{Mor} %Morphismen

% \DeclareMathOperator{\TOP}{TOP} %Topologische Räume
\newcommand{\TOP}{\textsc{Top}}
% \DeclareMathOperator{\HTOP}{HTOP} %Topologische Räume + Homotopieäquivalenz
\newcommand{\HTOP}{\textsc{HTop}}
% \DeclareMathOperator{\VR}{VR} %Vektorräume
\newcommand{\VR}{\textsc{VR}}
% \DeclareMathOperator{\MOD}{MOD} %Moduln
\newcommand{\MOD}{\textsc{Mod}}
% \DeclareMathOperator{\MENGEN}{MENGEN} %Mengen
\newcommand{\MENGEN}{\textsc{Mengen}}
% \DeclareMathOperator{\MAN}{MAN} %Mannigfaltigkeiten
\newcommand{\MAN}{\textsc{Man}}
% \DeclareMathOperator{\GRUPPEN}{GRUPPEN} %Gruppen
\newcommand{\GRUPPEN}{\textsc{Gruppen}}
% \DeclareMathOperator{\ABELGRUPPEN}{ABEL.GRUPPEN} %abelsche Gruppen
\newcommand{\ABELGRUPPEN}{\textsc{Abel.Gruppen}}
% \DeclareMathOperator{\KAT}{KAT} %Kategorien
\newcommand{\KAT}{\textsc{Kat}}
% \DeclareMathOperator{\FUN}{FUN} %Funktoren
\newcommand{\FUN}{\textsc{Fun}}


\newcommand{\op}{\mathrm{op}}

\newcommand{\Ker}{\ensuremath{\Kern \,}}
\newcommand{\zirkel}{\rotatebox[origin = c]{-90}{$\varangle$}}
\newcommand{\wirkung}{\!\curvearrowright\!}

\newcommand{\oE}{\OE}

%--Beweisende
\newcommand{\bewende}{\ifmmode \tag*{$\square$} \else \hfill $\square$ \fi}

%--Volumen
\newcommand{\vol}[1]{\ensuremath{\Vol_{#1}}}

%--Skalarprodukt
\DeclarePairedDelimiterX\sprod[2]{\langle}{\rangle}{#1\,\delimsize\vert\,#2}

%--Betrag, Gaußklammer
\DeclarePairedDelimiter{\abs}{\lvert}{\rvert}
\DeclarePairedDelimiter{\ceil}{\lceil}{\rceil}

%--Norm
\DeclarePairedDelimiter\doppelstrich{\Vert}{\Vert}
\newcommand{\norm}[2][\relax]{
\ifx#1\relax \ensuremath{\doppelstrich*{#2}}
\else \ensuremath{\doppelstrich*{#2}_{#1}}
\fi}


%--Umklammern
\DeclarePairedDelimiter\enbrace{(}{)}
\DeclarePairedDelimiter\benbrace{[}{]}
\DeclarePairedDelimiter\lenbrace{<}{>}

%--Mengen
\DeclarePairedDelimiterX\mengenA[1]{\lbrace}{\rbrace}{#1}
\DeclarePairedDelimiterX\mengenB[2]{\lbrace}{\rbrace}{#1\, \delimsize\vert \, #2}

\makeatletter
\newcommand{\set}[2][\relax]{
\ifx#1\relax \ensuremath{
\mengenA*{#2}}
\else \ensuremath{%
  \mengenB*{#1}{#2}}
\fi}
\makeatother


%--offen und abgeschlossen
\newcommand{\off}{\! \stackrel[\text{offen}]{}{\subset} \!}
\newcommand{\abg}{\! \stackrel[\text{abg.}]{}{\subset} \!}

%--Differential
\newcommand{\diff}[2]{\ensuremath{\frac{{\partial #1}}{{\partial #2}} }}
\newcommand{\diffd}[2]{\ensuremath{\frac{\mathd #1}{\mathd #2} }}

%--Diagonale Linien für Matrizen
\newcommand\tikzmark[1]{%
  \tikzexternaldisable\tikz[overlay,remember picture,baseline] \node [anchor=base] (#1) {};\tikzexternalenable}

\newcommand\MyLine[3][]{%
  \tikzexternaldisable\begin{tikzpicture}[overlay,remember picture]
    \draw[#1] (#2.south) -- (#3.north);
  \end{tikzpicture}\tikzexternalenable}
  
%--Jordankasten
\newcommand{\Jord}{ \tikzexternaldisable \begin{BMAT}{|ccc|}{|ccc|} %
				0 & & \\ %
				1\tikzmark{z} &  \diagdown & \\ %
				& \tikzmark{u}1 \MyLine[thick]{z}{u} &  0 %
			\end{BMAT} \tikzexternalenable}
			
%--Jordankasten mit Lambda
\newcommand{\JordLa}[1][\empty]{%
	\tikzexternaldisable\ifthenelse{\equal{#1}{\empty}}
		{\begin{BMAT}{|ccc|}{|ccc|} %
				\lambda & & \\ %
				1\tikzmark{z} &  \diagdown & \\ %
				& \tikzmark{u}1 \MyLine[thick]{z}{u} &  \lambda  %
			\end{BMAT}}
		{\begin{BMAT}{|ccc|}{|ccc|} %
				\lambda_{#1}  & & \\ %
				1\tikzmark{z} &  \diagdown & \\ %
				& \tikzmark{u}1 \MyLine[thick]{z}{u} &  \lambda_{#1}  %
			\end{BMAT}}\tikzexternalenable}

%--Inhaltsverzeichnis
\usepackage[tocindentauto]{tocstyle}
\usetocstyle{KOMAlike}			

\newcommand{\fach}{Differentialgeometrie \RM{1}.}
\newcommand{\semester}{WiSe 2014}
\newcommand{\homepage}{https://wwwmath.uni-muenster.de/u/frederik.witt/diffgeoI.html}
\newcommand{\prof}{Frederik Witt}
\newcommand{\verfasser}{Arne Grauer}
\providecommand{\verfasser}{Jannes Bantje}

%--Konfiguration von scrheadings
\setheadsepline{1pt}[\color{light_gray}]
\pagestyle{scrheadings}
%\fancyhf[H, F]{}
\clearscrheadfoot

\providecommand{\shortFach}{\fach}
\lehead{ \includegraphics[height=0.5 cm,keepaspectratio]{../!config/Bilder/Logo_WWU_Muenster_light_gray.pdf}}
\rehead{\rule{0cm}{0.5cm}\footnotesize \sffamily \color{light_gray} \verfasser -- Mitschrift \shortFach}
\lohead{\rule{0cm}{0.5cm} \footnotesize \sffamily \color{light_gray} Stand: \today \; \thistime[:]}
\rohead{\includegraphics[height=0.5 cm,keepaspectratio]{../!config/Bilder/fb10logo_gray.pdf}}


\ofoot[{ \color{dark_gray} \LARGE \sffamily \thepage}]{{ \color{dark_gray} \LARGE \sffamily \thepage}} %hier wir auch der plain Stil bearbeitet!
\automark{section}
\ifoot{ \color{dark_gray} \small \leftmark}

%--Metadaten
\author{\verfasser}
\titlehead{\includegraphics[height=1.5cm, keepaspectratio]{../!config/Bilder/Logo_WWU_Muenster.pdf}%
\hfill \includegraphics[height=1.3cm, keepaspectratio]{../!config/Bilder/fb10logo.pdf}}
\title{Skript \fach}
\subtitle{Mitschrift der Vorlesung  \enquote{\fach} von \prof}

\publishers{
Aktuelle Version verfügbar bei: \bigskip\\
\normalsize
\begin{minipage}[t]{0.45\textwidth}
	\centering{
	\qrcode[height=3.5cm, version=4]{https://github.com/JaMeZ-B/latex-wwu}
	\medskip\\
	\includegraphics[height=0.4cm, keepaspectratio]{../!config/Bilder/github_octo.pdf} 
	\includegraphics[height=0.4cm, keepaspectratio]{../!config/Bilder/GitHub_Logo.pdf} (inklusive Sourcecode)\\
	{\small\url{https://github.com/JaMeZ-B/latex-wwu}}}
\end{minipage}
\hfill
\begin{minipage}[t]{0.45\textwidth}
	\centering{
	\qrcode[height=3.5cm, version=4]{btsync://B6WH2DISQ5QVYIRYIEZSF4ZR2IDVKPN3I?n=latex_share}
	\medskip\\
	\includegraphics[height=0.4cm, keepaspectratio]{../!config/Bilder/bt_sync_logo.pdf} 
	{\large \textbf{Bittorrent} Sync} \\
	\texttt{B6WH2DISQ5QVYIRYIEZSF4ZR2IDVKPN3I}
	}
\end{minipage}
}



\begin{document}
\pagenumbering{Roman}
\maketitle
\begin{abstract}
\section*{Vorwort --- Mitarbeit am Skript}
Dieses Dokument ist eine Mitschrift aus der Vorlesung \enquote{\fach, \semester}, gelesen von \prof. Der Inhalt entspricht weitestgehend dem Tafelanschrieb. Für die
Korrektheit des Inhalts übernehme ich keinerlei Garantie! Für Bemerkungen und Korrekturen -- und seien es nur Rechtschreibfehler -- bin ich sehr dankbar. 
Korrekturen lassen sich prinzipiell auf drei Wegen einreichen: 
\begin{itemize}
	\item Persönliches Ansprechen in der Uni, Mails an \hrefsymmail{mailto:\mail}{\mail} (gerne auch mit annotieren PDFs) 
	\item \emph{Direktes} Mitarbeiten am Skript: Den Quellcode poste ich auf GitHub (siehe oben), also stehen vielfältige Möglichkeiten der Zusammenarbeit zur Verfügung:
	Zum Beispiel durch Kommentare am Code über die Website und die Kombination Fork + Pull Request. Wer sich verdient macht oder ein Skript zu einer Vorlesung, die 
	ich nicht besuche, beisteuern will, dem gewähre ich gerne auch Schreibzugriff.
	
	Beachten sollte man dabei, dass dazu ein Account bei \url{github.com} notwendig ist, der allerdings ohne Angabe von persönlichen Daten angelegt werden kann. 
	Wer bei GitHub (bzw. dem zugrunde liegenden Open-Source-Programm \enquote{\texttt{git}}) -- verständlicherweise -- Hilfe beim Einstieg braucht, dem helfe ich gerne 
	weiter. Es gibt aber auch zahlreiche empfehlenswerte Tutorials im Internet.\footnote{zB. \url{https://try.github.io/levels/1/challenges/1}, ist auf Englisch, aber dafür 
	interaktives LearningByDoing}
	\item \emph{Indirektes} Mitarbeiten: \TeX-Dateien per Mail verschicken. 
	
	Dies ist nur dann sinnvoll, wenn man einen ganzen Abschnitt ändern möchte (zB. einen alternativen Beweis geben), da ich die Änderungen dann per Hand einbauen muss! Ich freue mich aber auch über solche Beiträge!
\end{itemize}
\section*{Vorlesungshomepage}
{\centering 
\begin{minipage}[c][][c]{\textwidth}
	\centering \qrcode[height=3cm]{\homepage} \medskip\\
	\footnotesize \url{\homepage}
\end{minipage}
\par}

\section*{Über die Differentialgeometrie}
\begin{description}
	\item[1. Differentialgeometrie] Mannigfaltigkeiten $\Rightarrow $ Differentialtopologie
	\item[2. Riemannsche Geometrie] Längen und Winkel messen $\leadsto$ Geodäte (lokal) kürzeste Verbindung zwischen 2 Punkten einer Mannigfaltigkeit.
	
	Krümmung
	
	Allgemeine Relativitätstheorie: Krümmung = Gravitation ist für solche Effekte verantwortlich
\end{description}
\end{abstract}

\setcounter{page}{1}
\tableofcontents
\cleardoubleoddemptypage
\pagenumbering{arabic}
\setcounter{page}{1}


\section{Differenzierbare Mannigfaltigkeiten} % (fold)
\label{sec:1}
Konvention: glatt = $C^\infty$.\\
Wiederholung: Diffeomorphismus $f : U \subseteq \mathds{R}^n \to \mathds{R}^n$ glatt, bijektiv und $f ^{-1}$ glatt.

\subsection{Untermannigfaltigkeiten} % (fold)
\label{sub:11}

\subsubsection[Definition: Submersion, Immersion und Étale]{Definition} % (fold)
\label{ssub:111}
Eine Abbildung $f : U \subseteq \mathds{R}^n \to \mathds{R}^m$ glatt heißt \Index{Submersion}/\Index{Immersion}/\Index{Étale} in $x \in U$, falls
$\mathd_x f : \mathds{R}^n \to \mathds{R}^m$ surjektiv/injektiv/Isomorphismus ist. \marginnote{Insbesondere $n\ge m$/$n \le m$/$n=m$}. $f$ heißt schlechthin
\bet{Submersion}/\bet{Immersion}/\bet{Étale}, falls $f$ Submersion/Immersion/Étale für alle $x \in U$ ist.
% subsubsection 111 (end)

\subsubsection[Bemerkung: Submersiv/immersiv/étale sind offene Bedingungen]{Bemerkung} % (fold)
\label{ssub:112}
Submersiv/immersiv/étale sind offene Bedingungen, d.h. ist $f$ Submersion/Immersion/Étale in $x$, so auch in einer hinreichend kleinen Umgebung von $x$ 
(vergleiche auch Satz 15 unten)
% subsubsection 112 (end)

\subsubsection{Satz über inverse Funktionen (inverse function theorem)} % (fold)
\label{ssub:113}
Sei $f : U \subseteq \mathds{R}^n \to \mathds{R}^m$ étale in $x$. Dann existiert eine offene Menge $x \in V \subset U \subset \mathds{R}^n$, sodass $f\big|_{V}$ ein
Diffeomorphismus auf seinem Bild ist, d.h. $f\big|_{V} : V \xrightarrow{\sim} f(V) $ ist bijektiv und $f\big|_V$, $(f\big|_V)^{-1}$ glatt.
% subsubsection 113 (end)

\subsubsection{Satz über implizite Funktionen (implicit function theorem)} % (fold)
\label{ssub:114}
Sei $f : U \subset \mathds{R}^{n+k} \to \mathds{R}^n$ eine Submersion in $0$ mit $f(0)=0$. Dann existiert eine glatte Funktion $g$ auf einer Umgebung von $0 \in \mathds{R}^n$
mit Werten in $\mathds{R}^k$, d.h.
\[
	g=g(x_1, \ldots , x_n) = \begin{pmatrix}
		g_1(x_1, \ldots ,x_n) \\
		\vdots \\
		g_k(x_1, \ldots ,x_n)
	\end{pmatrix}
\]
mit
\[
	f \enbrace[\big]{x_1, \ldots ,x_n , g_1(x_1, \ldots , x_n), \ldots , g_k(x_1, \ldots ,x_n)} = f \enbrace[\big]{x,g(x)} = 0
\]
% subsubsection 114 (end)

\subsubsection{Satz (Normalenformen von Submersionen/Immersionen)} % (fold)
\label{ssub:115}
\begin{enumerate}[a)]
	\item Sei $f : U \subseteq \mathds{R}^{n+k} \to \mathds{R}^n$ eine Submersion mit $f(0)=0$. Dann existiert ein Diffeomorphismus 
	$\varphi : 0 \in V \subset \mathds{R}^{n+k} \to \mathds{R}^{n+k}$ mit $\varphi(V) \subset U$ und 
	$f \circ \varphi(x_1, \ldots , x_{n+k}) = (x_1, \ldots , x_n) = \pi_{\mathds{R}^n}(x_1, \ldots , x_{n+k})$, d.h. modulo einer Diffeomorphismus ist jede Submersion 
	äquivalent zu einer Projektion.
	\item Sei $f : U \subset \mathds{R}^n \to \mathds{R}^{n+k}$ eine Immersion in $0$, $f(0)=0$. Dann existiert ein Diffeomorphismus 
	\[
		\varphi : 0 \in V \subset \mathds{R}^n \to \mathds{R}^{n+k}
	\]
	mit $\varphi(0)=0$ und 
	\[
		\varphi \circ f (x_1, \ldots ,x _n) = \enbrace*{x_1, \ldots , x_n, 0 , \ldots , 0} = i(x_1, \ldots , x_n) 
	\]
	d.h. modulo Diffeomorphismen ist $f$ äquivalent zur kanonischen Einbettung von $\mathds{R}^n \hookrightarrow \mathds{R}^{n+k}$.
\end{enumerate}
\minisec{Beweis}
\begin{enumerate}[(i)]
	\item Wähle Koordinaten $x_1, \ldots , x_n$ mit $\mathd_0 f = (A \, B)$ mit $A \in \mathds{R}^{n \times n}$, $B \in \mathds{R}^{k \times n}$ mit $\det A \not= 0$.
	Definiere $F : U \subset \mathds{R}^n \times \mathds{R}^k \to \mathds{R}^n \times \mathds{R}^k $ durch $(p,q) \mapsto \enbrace*{f(p,q), q}$. Dann gilt
	\[
		\mathd_{(0,0)} F = \begin{pmatrix}
			\mathd_0 f & \\
			0 & \id_{k \times k}
		\end{pmatrix} = \begin{pmatrix}
			A & B \\
			0 & \id_{k \times k} 
		\end{pmatrix}
		\Longrightarrow \det \mathd_0 F \not= 0
	\]
	Mit \ref{ssub:113} folgt: $\varphi = F ^{-1}$ Diffeomorphismus um $0 \in \mathds{R}^{n+k}$ herum.
	\item Definiere $F : U \times \mathds{R}^k \to \mathds{R}^n \times \mathds{R}^k$. $F(p,q) = f(p)+ (0,q)$.
	\[
		\mathd_{(0,0)} F = \begin{pmatrix}
			\mathd_0 f & 0 \\
			0 & \id_{k \times k}
		\end{pmatrix},
	\]
	also $\det \mathd_{(0,0)} F \not= 0$. $\Rightarrow  \exists \varphi = F ^{-1}$ um $0$ mit $\varphi(f(p))=\varphi \circ F (p,0) = (p,0) = i(p)$ \bewende
\end{enumerate}
\minisec{Bemerkung}
Durch Komposition mit Translationen können Submersionen/Immersionen mit $F(x)=y$ (nicht notwendigerweise $=0$) behandelt werden.
% subsubsection 115 (end)

\subsubsection[Definition: Untermannigfaltigkeit]{Definition} % (fold)
\label{ssub:116}
$M \subset \mathds{R}^{n+k}$ heißt $n$-dimensionale \Index{Untermannigfaltigkeit}, falls $\forall x \in M : \exists x \in U \subset \mathds{R}^{n+k}$ Umgebung, sowie eine
Submersion $f : U \to \mathds{R}^k$ mit $U \cap M = f ^{-1}(0)$. \marginnote{Verallgemeinerung des $n$-dimensionalen Unterraumes}
% subsubsection 116 (end)

\subsubsection[Beispiele für Untermannigfaltigkeiten]{Beispiele} % (fold)
\label{ssub:117}
\begin{enumerate}[(i)]
	\item Die $n$-\Index{Sphäre} $S^n := \set[x=(x_1, \ldots ,x_{n+1}) \in \mathds{R}^{n+1}]{\abs*{x}^2=1 } $ ist eine $n$-dimensionale Untermannigfaltigkeit des 
	$\mathds{R}^{n+1}$. Definiere $f : \mathds{R}^{n+1} \to \mathds{R}$ durch $f(x_1, \ldots , x_{n+1}) = x_1^2 + \ldots + x_{n+1}^2 -1$. $\Rightarrow f ^{-1}(0) = S^n$ nach
	Definition.
	\[
		\mathd_x f = (2 x_1, \ldots , 2 x_{n+1}) : \mathds{R}^{n+1} \to \mathds{R}
	\]
\end{enumerate}
% subsubsection 117 (end)
% subsection 11 (end)
% section 1 (end)


\cleardoubleoddemptypage
\pagenumbering{Alph}
\setcounter{page}{1}
\printindex
\listoffigures
\end{document}
